\documentclass[11pt,a4paper]{jlreq}
\usepackage{amsmath,amssymb}
\usepackage{booktabs}
\usepackage{array}

\begin{document}

\begin{center}
{\LARGE \textbf{情報科学基礎A 第4回 レポート}}
\end{center}
\vspace{1em}

\begin{flushright}
提出日：\today \\
学籍番号：2611193 \\
氏名：中井平蔵
\end{flushright}

\section*{Q1: \boldmath$H(X,Y)=H(Y)+H(X\mid Y)$\unboldmath の証明}

右辺を整理すると、

\begin{align*}
H(Y)+H(X\mid Y)
&=\sum_{y\in A_Y} p(y)\log\frac{1}{p(y)}
 +\sum_{y\in A_Y}p(y)\sum_{x\in A_X}p(x\mid y)\log\frac{1}{p(x\mid y)}\\
\end{align*}

ここで、
\[
\sum_{x\in A_X}p(x,y)=p(y)
\]

より、

\begin{align*}
H(Y)+H(X\mid Y)
&=\sum_{y\in A_Y}\sum_{x\in A_X}p(x,y)\log\frac{1}{p(y)}
 +\sum_{y\in A_Y}p(y)\sum_{x\in A_X}p(x\mid y)\log\frac{1}{p(x\mid y)}\\
&=\sum_{y\in A_Y}\sum_{x\in A_X} p(x,y)\log\frac{1}{p(y)}
 +\sum_{y\in A_Y}\sum_{x\in A_X} p(x,y)\log\frac{1}{p(x\mid y)}\\
&=\sum_{y\in A_Y}\sum_{x\in A_X} p(x,y)\left(\log\frac{1}{p(y)}+\log\frac{1}{p(x\mid y)}\right)\\
&=\sum_{y\in A_Y}\sum_{x\in A_X} p(x,y)\log\frac{1}{p(y)p(x\mid y)}\\
&=\sum_{y\in A_Y}\sum_{x\in A_X} p(x,y)\log\frac{1}{p(y)\cdot\frac{p(x,y)}{p(y)}}\\
&=\sum_{y\in A_Y}\sum_{x\in A_X} p(x,y)\log\frac{1}{p(x,y)}\\
&=H(X,Y)
\end{align*}

したがって、
\[
H(X,Y)=H(Y)+H(X\mid Y)
\]

が成り立つ。

\section*{Q2: 128ビット列の作成とエントロピー計算}

以下はランダムに生成した128ビット列である。

\[
\texttt{1100101011110001010100011101011010001010101110010101010001110010}
\]
\[
\texttt{0010100110101110101001000101110101010001011101000101111110001010}
\]

\(N(x)\) を、記号 \(x\) の度数とする。

\subsection*{(1) \boldmath$\{0,1\}^{128}$ とみなす場合}

$0$ と $1$ の度数は、
\[
N(0)=64,\qquad N(1)=64
\]

より、
\[
\hat p(0)=\frac{64}{128}=0.5000,\qquad
\hat p(1)=\frac{64}{128}=0.5000
\]

である。

したがって、1ビットあたりのエントロピーは、
\begin{align*}
H_1
&=-\sum_{b\in\{0,1\}}\hat p(b)\log_2\hat p(b)\\
&= -\frac{64}{128}\log_2\frac{64}{128}
   -\frac{64}{128}\log_2\frac{64}{128}\\
&=1.0000\ \text{bit/記号}
\end{align*}

となる。

\subsection*{(2) \boldmath$\{00,01,10,11\}^{64}$ とみなす場合}

度数は次の通りである。
\[
N(00)=11,\ \ N(01)=23,\ \ N(10)=19,\ \ N(11)=11
\]

よって、
\[
\hat p(00)=\frac{11}{64},\ \hat p(01)=\frac{23}{64},\ \hat p(10)=\frac{19}{64},\ \hat p(11)=\frac{11}{64}
\]

であり、エントロピーは、
\begin{align*}
H_2
&=-\sum_{u\in\{00,01,10,11\}}\hat p(u)\log_2\hat p(u)\\
&\approx 1.9241\ \text{bit/2ビット記号}
\end{align*}

となる。

比較のため1ビットあたりに直すと、
\[
\frac{H_2}{2}\approx 0.9620\ \text{bit/bit}
\]

となる。

\subsection*{(3) \boldmath$\{0000,0001,\dots,1111\}^{32}$ とみなす場合}

度数は次の通り。
\begin{table}[h]
\centering
\caption{4ビット記号ごとの度数}
\begin{tabular*}{0.45\linewidth}{@{\extracolsep{\fill}}>{\centering\arraybackslash}p{0.2\linewidth}>{\centering\arraybackslash}p{0.2\linewidth}}
\toprule
$v$ & $N(v)$\\
\midrule
0000 & 0\\
0001 & 3\\
0010 & 2\\
0011 & 0\\
0100 & 3\\
0101 & 5\\
0110 & 1\\
0111 & 2\\
1000 & 2\\
1001 & 2\\
1010 & 5\\
1011 & 1\\
1100 & 1\\
1101 & 2\\
1110 & 1\\
1111 & 2\\
\bottomrule
\end{tabular*}
\end{table}

したがって、
\[
H_4=-\sum_{v\in\{0000,\dots,1111\}}\hat p(v)\log_2\hat p(v)
\approx 3.6022\ \text{bit/4ビット記号}
\]

となる。

1ビットあたりに直すと、
\[
\frac{H_4}{4}\approx 0.9006\ \text{bit/bit}
\]

となる。

\subsection*{考察}

今回の結果は、
\[
H_1=1.0000,\qquad \frac{H_2}{2}\approx 0.9620,\qquad \frac{H_4}{4}\approx 0.9006
\]
となった。$1$ビット単位の値 \(H_1=1.0000\) は最大値に一致しており、
ランダム性が十分に高い。

一方で、ブロック長を$2$ビット、$4$ビットと長くすると、$1$ビットあたりのエントロピーは、
\[
1.0000 \to 0.9620 \to 0.9006
\]

と小さくなった。これは、ブロック長を長くするほど並び方の種類が増えるのに対し、
観測できるブロック数が減るためである。
今回の長さは$128$ビットであるため、$4$ビットに分けると$32$個しか数えられず、$16$種類ある並びの中には$1$回も現れないものも出る。
よって、一様分布から外れやすくなり、エントロピーが小さくなる。

\end{document}
